\section{System Overview}\label{sec:overview}

\begin{figure}[t]
\centering
\begin{tikzpicture}[
  font=\footnotesize,
  >={Stealth[length=4pt,width=3.5pt]},
  comp/.style={draw, rounded corners=1.5pt, align=center, inner sep=2pt,
    minimum height=0.8cm, text width=#1},
  comp/.default=2.1cm,
  trusted/.style={fill=black!12},
  keyedge/.style={->, line width=1.3pt},
  aux/.style={->, black!55},
  bnd/.style={draw, dashed, rounded corners=3pt, inner sep=2.5pt},
]
% row 0: outside the key path
\node[comp=1.6cm] (br) at (0.85,0) {Browser};
\node[comp]       (po) at (3.95,0) {Portals (Django)\\no subject keys};
\node[comp]       (st) at (7.3,0) {Postgres store\\ciphertext, wrapped\\keys, index, log};
% row 1: key holders
\node[comp=1.6cm] (an) at (0.85,-1.6) {Authority net\\share $x{=}2$};
\node[comp, trusted] (cu) at (3.95,-1.6) {Key Custodian\\all persona keys\\share $x{=}1$};
\node[comp, trusted] (au) at (7.3,-1.6) {Audit service\\hash chain\\shares $x{=}3,4$};
% row 2
\node[comp=1.6cm] (am) at (0.85,-3.0) {Authority med};
\node[comp]       (fa) at (3.95,-3.0) {Facility A\\threshold gate};
\node[comp]       (fb) at (7.3,-3.0) {Facility B\\threshold gate};
% trust boundaries
\begin{scope}[on background layer]
  \node[bnd, fit=(po)(st), label={[yshift=-1pt]above:{store data, identify users, decide no access}}] {};
  \node[bnd, fit=(an)(am)] {};
  \node[bnd, fit=(fa)(fb), label={[yshift=1pt]below:{facilities: hold no records}}] {};
  \node[bnd, fit=(cu)] {};
  \node[bnd, fit=(au)] {};
\end{scope}
% non-key traffic
\draw[aux] (br) -- node[above] {HTTP} (po);
\draw[aux] (po) -- node[above] {read} (st);
\draw[aux] (po) -- (cu);
\draw[aux] (fb) -- (au);
% key path
\draw[keyedge] (an) -- node[above] {$x{=}2$} (cu);
\draw[keyedge] (au) -- node[above] {$x{=}3,4$} (cu);
\draw[keyedge] (cu) -- (fa);
\draw[keyedge] (cu.south east) -- (fb.north west);
\end{tikzpicture}
\caption{Simulated federation; dashed boxes are trust boundaries, shaded
parties are trusted. Thick arrows carry emergency-key shares, sealed to a
facility's gate; share~2 is used only if share~1 is missing. Thin arrows:
portal to custodian by subject id, no keys; facility to audit, log append
and, at rotation, the new sealed shares~3,4. Not drawn: portal-to-facility
calls, rotation deposits to custodian and authority, Postgres, root CA.}
% src: facts-F4#43, facts-F2#85, facts-F4#10, facts-F4#12, facts-F4#20, facts-F4#22, facts-F5#9, facts-F5#10
\label{fig:components}
\end{figure}

The federation is simulated: two mock facilities, two attribute authorities,
an audit service, a Key Custodian, a PostgreSQL store and one Django process
serving the patient, clinician and admin portals run on one machine, as local
processes or as Docker Compose containers on one bridge, under one
self-signed PKI root that signs both facility CAs
(Fig.~\ref{fig:components}). % src: facts-F4#1, facts-F4#5, facts-F4#27, facts-F4#43
There is no health-information exchange and no FHIR/HL7 interface, and all records are
synthetic. % src: facts-F4#1, facts-F4#44, facts-F4#39
The emergency scenario assumes the patient can be identified
(Section~\ref{sec:limits}, item~7). % src: facts-F9#7

\textbf{Components.} Facility~A is the patient's home facility and Facility~B
the out-of-network one where the emergency happens; neither holds records, and
both read ciphertext from the shared store and run the policy engine and
threshold gate. % src: facts-F4#7
An \emph{attribute authority} issues credentials and attribute keys, and the
patient's \emph{home authority} is the one that issued the patient's
\texttt{owner} credential. % src: facts-F4#8, facts-F5#12
The \emph{Key Custodian} is the key agent of clinicians and patients.
% src: facts-F4#14
This build chooses a weighted 3-of-4 split of each emergency key: $x{=}1$ to the custodian, $x{=}2$ to the home authority
(always \texttt{authority-net} here), $x{=}3,4$ to
the audit service, so every set of three shares includes an audit
share. % src: facts-F4#10, facts-F4#11, facts-F2#36

\textbf{Trust boundaries.} Every service call carries a root-issued service
certificate and a signed one-time request envelope, checked before any route
code runs, and key-bearing facility routes admit only the custodian or
bootstrap. % src: facts-F4#16, facts-F4#15
Shares travel sealed to the facility's gate; the audit service releases its
shares only after validating the assertion and committing a log entry
(Section~\ref{sec:breakglass}). % src: facts-F4#12

\textbf{Identity and certificates.} One self-signed root signs every facility
CA, attribute authority and service certificate; its key is held only at
bootstrap and is a single point of trust. % src: facts-F4#27, facts-F4#31
Certificates are X.509 with ECDSA P-256; a clinician's is signed by their
facility CA. % src: facts-F4#28, facts-F4#29
Because both facility CAs chain to the same root, Facility~B can validate a
clinician enrolled at Facility~A, which cross-facility break-glass needs;
chain validation fails closed on any bad link. % src: facts-F4#27, facts-F4#30
Each attribute authority issues signed credentials in its own namespace and
releases an attribute key only against a valid, unrevoked, unexpired
credential. % src: facts-F4#8, facts-F4#32

\textbf{Discovery and identification.} Records live in one shared store, so
discovery is a patient-indexed lookup that returns record identifiers only,
not a search across facilities. % src: facts-F4#9
To find the patient, the clinician portal looks up a name, optionally with a
date of birth, in its own directory; ambiguous matches are listed and never
auto-selected, a lookup that cannot be audited is not shown, and identifying
a patient releases no key. % src: facts-F9#7, facts-F4#19

\textbf{What portals and Postgres may not do.} They store data and identify
users but never decide access. % src: facts-F4#19
The portal holds no subject key and sends every key-needing action to the
custodian, so a portal account alone decrypts nothing; its database roles
cannot write store tables, and its consent state is display data only. % src: facts-F4#20, facts-F4#21, facts-F4#22, facts-F4#23
